\subsection{Excluding a point of order $77$}
\begin{lemma}\label{lem:decic77-finite-image-Cartier}
Let $Y/\mathbb Q$ be an odd-level modular curve or diamond quotient,
with good reduction at $2$, finite etale cyclotomic cuspidal
boundary, and a rational cusp $c_0$.
Suppose that $T_2$ and $\delta=\langle2\rangle$ descend to its
Jacobian $J$, satisfy Eichler--Shimura, and $T_2J(\mathbb Q)$ is
finite. Put
\[
 U=\mathcal C(\mathcal C-1)H^0(Y_{\mathbb F_2},\Omega^1).
\]
If restriction from $U$ has rank $d$ at an effective cuspidal
divisor $\overline C$ of degree $d$, no rational noncuspidal
effective divisor of degree $d$ on $Y$ reduces to $\overline C$.
\end{lemma}
\begin{proof}
Let $C$ be the etale cuspidal lift of $\overline C$.
Each closed residue cusp is lifted by its unramified local trace,
so $C$ is a divisor over $\mathbb Q_2$ even when its individual
geometric cusps are not rational over $\mathbb Q_2$.
Suppose that $D$ is a proposed rational divisor.
The class
\[
 T_2[D-C]=T_2[D-dc_0]-T_2[C-dc_0]
\]
is torsion by the finite-image hypothesis and Manin--Drinfeld.
It belongs to $J(\mathbb Q_2)$ and reduces to zero.
Lemma~\ref{lem:dyadic-formal-Hecke-kernel} gives
\[
 (T_2-\delta)T_2[D-C]=0.
\]
The morphism from $Y^{(d)}$ to the full Jacobian, given by
$E\mapsto(T_2-\delta)T_2[E-dc_0]$, therefore has the same value
at $D$ and $C$.
On the special fibre its cotangent image is
\[
 (\delta^*)^2\mathcal C(\mathcal C-1)H^0(\Omega^1)=U.
\]
Here Frobenius has zero differential, Verschiebung pulls back
as Cartier, and the invertible map $\delta^*$ commutes with Cartier.
The restriction condition gives formal immersion at $\overline C$.
Thus the two local sections coincide, contradicting the
noncuspidal support of $D$.
\end{proof}

\begin{lemma}\label{lem:decic77-geometry}
Put $X=X_1(77)$ and
$Y=X/\langle\langle36\rangle\rangle$.
Then $Y$ has good reduction at $2$, genus $37$, and
$T_2J_Y(\mathbb Q)$ is finite.
The cover $Y\to X_0(77)$ is etale cyclic of degree $6$.
Its deck group is generated by $\delta=\langle2\rangle$, and
$\langle3\rangle=\delta^{-1}$.
\end{lemma}
\begin{proof}
Let $\pi:X\to Y$ be the quotient by
$\langle\langle36\rangle\rangle$, of order $5$.
This group acts freely, including on the special fibre at $2$.
Thus $Y$ has good reduction at $2$, and
$g(Y)=1+(181-1)/5=37$.
The complete degree-independent winding computation accounts
for all $181$ source dimensions.
Every nontrivial-character component is full winding.
The trivial-character component has dimension $7$ and winding
dimension $6$, and $T_2$ is zero on its one-dimensional quotient.
Kato's theorem, in the winding formulation of
\cite[Remarks~3.2 and 3.3]{DEHMZ}, therefore gives that $T_2J_X(\mathbb Q)$ is finite.
Pullback and norm show that $T_2J_Y(\mathbb Q)$ is finite as well.

We retain the exact modular identification of $Y_{\mathbb F_2}$.
The base $X_0(77)$ is the genus-$7$ normalization obtained from
$X_0(11)$ by adjoining $v$ with
\[
 v^8+jv^7+v^4+1=0.
\]
Its Picard group is $C_{30}\times C_{60}\times\mathbb Z$.
The degree-$6$ cyclic cover has cusp Frobenius elements
$0,2,3,5$ at widths $77,11,7,1$.
The cusp conditions leave six characters to $C_6$; splitting
at two actual degree-$10$ source images singles out the modular
character.
Its cubic subcover is
$X_1(7)\times_{X(1)}X_0(11)$, with Tate parameter $t$ and
\[
 b=t^3-t^2,\quad c=t^2-t,\quad
 v=\frac{t^2+t}{t^3+t+1}.
\]
The fixed Artin--Schreier component is checked on generators of
the base Picard group modulo $2$.
The composite has genus $37$, a rational cusp and $j$-degree
$576$. Its two physical deck transformations are rebuilt and
verified exactly.


Since $2^6=36^2$ and no smaller positive power of $2$ belongs
to $\{\pm36^i:0\leq i<5\}$, $\delta$ has order $6$.
Also $36^3=71\equiv-6\pmod{77}$, giving
$\langle3\rangle\delta=1$.
In the explicit cubic and Artin--Schreier model, with coordinates
$(t,s,z)$, its physical action is
\[
 \delta(t,s,z)=\left(\frac{t+1}{t},s,z+1\right).
\]
The Tate normal form of $2P_7$ verifies the first coordinate;
the quadratic character has value $-1$ at $2$.
Both component actions and the order-$6$ identity are checked
exactly on the defining function fields.
\end{proof}

\begin{lemma}\label{lem:decic77-pure-cusps}
A rational effective divisor of degree $10$ on $X_1(77)$ with
entirely noncuspidal support cannot have entirely cuspidal
reduction at $2$.
\end{lemma}
\begin{proof}
Suppose that $D$ is such a divisor with entirely cuspidal reduction.
Use $\pi:X\to Y$ from Lemma~\ref{lem:decic77-geometry}.
The source cusp signature is $1^{30},3^{10},10^3,30$.
On $Y$, the width-degree pairs are
\[
 (77,1)^6,\qquad(11,3)^2,\qquad(7,2)^3,\qquad(1,6).
\]
A degree-$10$ source cusp divisor therefore pushes forward
either to a divisor on the six rational and two cubic cusps,
or to $5Q$ for one of the three quadratic cusps.
The latter weight is the residue degree ratio $10/2$.
Exact primitive-column orbits verify this complete coverage.
The first kind gives
\[
 \sum_{r=0}^3(r+1)\binom{15-3r}{5}=4989
\]
patterns; the three additional images give $4992$ in total.

For $U=\mathcal C(\mathcal C-1)H^0(Y_{\mathbb F_2},\Omega^1)$,
we have $\dim U=32$.
In one exact differential basis the $69$ restriction blocks
cover rational multiplicities through $10$, cubic multiplicities
through $3$, and multiplicity $5$ at each quadratic cusp.
All $4992$ restrictions have rank $10$.
The rational multiplicity-$10$ and quadratic multiplicity-$5$
blocks are freshly computed; verified deck transformations
transport them to their six and three cusp orbits.
An independent binary replay identifies all $60$ retained
lower-multiplicity kernels across the exact basis transition,
and verifies every restriction rank.
A further native calculation in a fresh model independently
gives rank $10$ at $10P$, for one rational cusp $P$, and at
$5Q$, for one quadratic cusp $Q$.
Deck symmetry covers the six and three corresponding patterns.
Applying Lemma~\ref{lem:decic77-finite-image-Cartier} to
$\pi_*D$ gives the contradiction.
\end{proof}


\begin{proposition}\label{prop:decic-no77}
An elliptic curve over a number field of degree $10$ has no point
of order $77$.
\end{proposition}
\begin{proof}
Put $X=X_1(77)$, and use the quotient $\pi:X\to Y$ from
Lemma~\ref{lem:decic77-geometry}.
A proposed torsion point gives a rational effective noncuspidal
divisor $D$ of degree $10$ on $X$.
The complete noncuspidal closed-point counts through degree $10$
at $2$ are
\[
 (0,0,0,0,0,0,0,0,0,90).
\]
An independent census checks every ordinary elliptic curve and
its quadratic twists; the supersingular trace possibilities
contribute no point of order $77$.
The $30$ qualifying curves over $\mathbb F_{1024}$ have cyclic
group of order $1078$ and form three $j$-orbits of degree $10$.
Each has a unique subgroup of order $77$, whose $30$ generator
classes form one source diamond orbit.
Thus the $90$ closed places form exactly three diamond orbits.

Lemma~\ref{lem:decic77-pure-cusps} excludes entirely cuspidal
reduction.
Any other reduction is one closed noncuspidal place of degree
$10$. Its image $P$ on $Y$ also has degree $10$, since its
$j$-invariant has that degree.
In particular there is no residue norm factor in this pushforward.

Put $\delta=\langle2\rangle$ and
$A_3=T_3-1-3\langle3\rangle$.
By Eichler--Shimura and injectivity of torsion reduction at the
good odd prime $3$, $A_3$ kills the rational torsion of $J_Y$.
Since $T_2J_Y(\mathbb Q)$ is finite, $A_3T_2$ kills
$J_Y(\mathbb Q)$.
On $J_Y(\mathbb F_2)$, Eichler--Shimura gives
$T_2=1+2\delta$.
Normalize at a rational cusp $c_0$.
The cusp formula gives $A_3c_0=0$ as a literal divisor, and also
$A_3\delta c_0=0$. At $2$ the same formula gives
$T_2c_0=c_0+2\delta c_0$, so $A_3T_2c_0=0$ as well.
Consequently the reduction of $\pi_*D$ would satisfy
\[
 [A_3(1+2\delta)P]=0.
\]
Using $\langle3\rangle=\delta^{-1}$, this is the principality
condition for
\begin{equation}\label{eq:decic77-B}
 B=T+2\delta T-7P-2\delta P-3\delta^{-1}P,
 \qquad T=T_3P.
\end{equation}
The positive and negative parts both have degree $120$ before
any cancellations.

We compute the required Hecke images on the explicit genus-$37$
model.
The marked point of order $7$ determines the Tate parameter $t$.
The cyclic subgroup of order $11$ determines its point on
$X_0(11)$, using the exact $j$-map and the $j$-invariant of its
$11$-isogeny quotient.
These coordinates give the image on the cubic subcover.
Both quadratic lifts have the same residue degree.
We fix either source lift and allow every quadratic lift of each
Hecke image. The deck involution covers the other source lift.
The complete tests are
\[
\begin{array}{c|c|c|c}
\text{case}&\text{kernel orbit lengths}&\text{lift allocations}
 &\text{nonprincipal tests}\\\hline
1&(1,3)&4&4\\
2&(1,1,1,1)&16&16\\
3&(1,3)&4&4.
\end{array}
\]
The image degrees are $10$ and $30$ in the first and third
cases, and $10$ in the second.
Every residue norm weight is one: $10w/d=1$ for each kernel
orbit of length $w$ and image degree $d$.
Thus these $24$ allocations contain the actual Hecke divisor
in every case.

For every allocation, the exact calculation gives $h^0(B)=0$,
contradicting \eqref{eq:decic77-B}.
The finite checks independently reconstruct all $12$ geometric
cyclic $3$-kernels and the $15$ marked cyclic-$11$ quotient maps
at the source and isogeny images.
They also verify the complete $3$-division polynomials, the marked
Tate coordinates, every norm weight and all quadratic splitting
conditions.
The operator divisors are exported before testing, using two
exact ideal generators for each prime and exact function
roundtrips and equality of the generated ideals.
A second native computation rebuilds the fields and prime ideals,
checks the source and diamond coordinates, and reproduces all
$24$ zero Riemann--Roch dimensions.
A separate native check binds all $24$ positive Hecke allocations
to the marked isogeny coordinates with their exact residue norms.
Hecke operators commute with the diamond group, so the three
case computations exclude their entire source diamond orbits.
This completes the proof.
\end{proof}
